\documentclass{beamer}
\usepackage{pgfpages}
\usepackage[LGR,T1]{fontenc}
\usepackage[utf8]{inputenc}
\usepackage[spanish]{babel}
\usepackage{textcomp}
\usepackage{alltt}
\renewcommand{\ttdefault}{txtt}
\usepackage{listings}
\lstset{language=c++, xleftmargin=25pt}
\input{listings-modelica.cfg}
\usepackage{amsmath}
\usepackage{mathtools}
\DeclarePairedDelimiter{\abs}{\lvert}{\rvert}
\usepackage{float}
\usepackage{graphicx}
\usepackage{caption}
\usepackage{subcaption}
\usepackage{epstopdf}

\setbeameroption{show notes on second screen}

\usetheme{Madrid}
 
\title[] %optional
{Causalización de Modelos Modelica}
 
\subtitle{Tesina de Grado}
 
\author[] % (optional, for multiple authors)
{Federico~Moya}
 
\institute[] % (optional)	
{
  Licenciatura en Cs. de la Computación\\
  Facultad de Cs. Exactas, Ingeniería y Agrimensura\\
  Universidad Nacional de Rosario
}
 
\date[] % (optional)
{Octubre 2016}
 
\logo{\includegraphics[height=1.5cm]{graphics/logo_unr.jpg}}
 
\AtBeginSection[]
{
  \begin{frame}
    \frametitle{Indice}
    \tableofcontents[currentsection]
  \end{frame}
}

\begin{document}

\begin{frame}
    \titlepage
\end{frame}

\section{Introducción}

\begin{frame}
	\frametitle{Contexto}
    \begin{itemize}
        \item<1-> Proyecto de investigación de la U.N.R: Modelado, Simulación y Control en Tiempo Real con Aplicaciones en Electrónica de Potencia.
        \item<2-> Desarrollar un compilador Modelica con los fines de investigar distintos algoritmos en relación a modelos de gran escala.
        \item<3-> Causalización de ecuaciones diferenciales algebraicas.
    \end{itemize}
\end{frame}

\section{Modelado y Simulación}

\begin{frame}
	\frametitle{Sistema}
	\begin{itemize}
		\begin{block}<1->{Definición}
			Un sistema es un objeto o una colección de objetos cuyas propiedades queremos estudiar.
		\end{block}
	  	\item<2-> Las entradas son variables del entorno que influencian el comportamiento del sistema. Estas entradas pueden o no ser controlables por nosotros. 
		\item<3-> Las salidas de un sistema son variables que son determinadas por el sistema y pueden influenciar el entorno que lo rodea.
	\end{itemize}
\end{frame}

\begin{frame}
	\frametitle{Experimento}
	\begin{block}<1->{Definición}
		Un experimento es el proceso de extraer información de un sistema ejercitando sus entradas.
	\end{block}
	Desventajas del método experimental:
	\begin{itemize}
	  	\item<2-> Para un gran número de sistemas muchas variables de entrada no son accesibles ni controlables.
		\item<3-> Algunas variables de salida importantes no resultan accesibles para poder medirlas; estas variables a veces se denominan estados internos del sistema.
		\item<4-> Problemas prácticos (el experimento puede ser muy caro, o peligroso, o bien el sistema necesario para el experimento puede que todavía no exista).
	\end{itemize}
\end{frame}

\begin{frame}
	\frametitle{Modelo}
	\begin{block}<1->{Definición}
		Un modelo de un sistema es cualquier cosa sobre la cual se puede aplicar un “experimento” con el objeto de responder preguntas sobre el sistema.
	\end{block}
	\begin{block}<2->{Modelo matemático}
		Es una descripción de un sistema en donde las relaciones entre las variables se expresan en forma matemática mediante ecuaciones. Las variables pueden ser cantidades medibles como tamaño, longitud, peso, temperatura, nivel de desempleo, flujo de información, etc. La mayoría de las leyes de físicas son modelos matemáticos en este sentido.
	\end{block}
\end{frame}

\begin{frame}
	\frametitle{Simulación}
	\begin{block}<1->{Definición}
		Una simulación es un experimento realizado sobre un modelo.
	\end{block}
    Algunos ejemplos:
	\begin{itemize}
  		\item<1-> La simulación de un proceso industrial como la producción de acero o papel para aprender sobre el comportamiento bajo diferentes condiciones de operación con el objetivo de mejorar el proceso.
        \item<2-> La simulación del comportamiento de un vehículo, por ejemplo un auto o un avión, para que el conductor o piloto pueda disponer de un entrenamiento realista.
        \item<3-> La simulación de un modelo simplificado de una red de computadoras para aprender cómo se comporta bajo diferentes cargas con el objetivo de mejorar la eficiencia.
  	\end{itemize}
\end{frame}


\section{Modelos Matemáticos para Sistemas Continuos}

\begin{frame}{Categorías de variables y constantes}
    \begin{itemize}
        \item<1-> \textbf{Constante:} una cantidad en el modelo que no varía nunca, por ejemplo la gravedad en la tierra $g=9.8$.
        \item<2-> \textbf{Variable:} una cantidad en el modelo que varia respecto del tiempo.
        \note{
            \begin{itemize}
                \item<1-> \textbf{Variable de salida:} una variable cuyo comportamiento queremos observar. Usualmente denotada como $y(t)$. El rol de variable de salida no está establecido por el sistema que se modela sino que es determinado por lo que el modelador quiere estudiar.
                \item<2-> \textbf{Variable de entrada:} una variable que afecta el comportamiento del modelo sin ser afectada por otras variables del modelo y cuyas variaciones respecto del tiempo pueden ser controladas. Usualmente se las denota $u(t)$.
            \end{itemize}
        }
        \item<3-> \textbf{Variable de estado}: una variable para la cual aparece su derivada respecto del tiempo, usualmente denotada por $\dot{x}(t)$
        \item<4-> \textbf{Variables algebraicas:} una variable que aparece en las ecuaciones pero que su derivada no aparece en el modelo.
    \end{itemize}
\end{frame}

\begin{frame}{Ecuacion diferencial ordinaria en forma de espacio de estados explícita}
    \note{ Las ecuaciones diferenciales ordinarias aparecen en la mayoría de los modelos matemáticos de sistemas continuos dado que la dependencia respecto del tiempo en esos modelos generalmente se refleja por medio de derivadas respecto del tiempo de variables de estado correspondientes a dichos modelos.\\
    Un modelo ODE de espacio de estados explícito resulta más eficiente de resolver y analizar que otros modelos más complicados.\\
    El modelo expresado de esta manera implica que a partir de las variables de estado y de las variables de entrada podemos calcular las derivadas y las variables algebraicas. }
    \framesubtitle{Forma General}
    \begin{align*}
        \dot{x}(t) & =f(x(t),u(t)) \\
        y(t)       & =g(x(t),u(t))
    \end{align*}
\end{frame}

\begin{frame}{Ecuacion diferencial ordinaria en forma de espacio de estados explícita}
    \framesubtitle{Ejemplo: Ley de fuerza de Newton}   
    \begin{center}
        $m\ddot{s}(t)=F(t)$
    \end{center}
    \pause 
    \begin{align*}
        \dot{v}(t) & =\frac{F(t)}{m}\\
        \dot{s}(t) & =v(t)
    \end{align*}
\end{frame} 

\begin{frame}[fragile]{Ecuacion diferencial algebraica}
   \note{
       Las ecuaciones diferenciales algebraicas no son los modelos más eficientes para resolver ya que las derivadas no siempre aparecen del lado izquierdo y por lo tanto puede ser necesario aplicar derivación numérica
   }
   \framesubtitle{Forma General}
   \begin{align*}
        F(x(t),\dot{x}(t),u(t),y(t)) & =0 \\
        G(x(t),u(t),y(t))            & =0
   \end{align*}
\end{frame}

\begin{frame}[fragile]{Ecuacion diferencial algebraica}
   \framesubtitle{Ejemplo: Circuito RLC} 
    \note{
        En el circuito hay cinco elementos, y cada uno define dos variables, denominadas la ensión y la intensidad de corriente a través de dicho elemento, necesitamos diez ecuaciones para describir el modelo, p. ej. las cinco ecuaciones correspondientes a cada elemento, las cuales definen la relación entre la tensión y la intensidad de corriente a través del elemento, más tres ecuaciones de malla en el voltaje de malla, más dos ecuaciones de nodos en la corriente del nodo.
    }
   % \fontsize{8pt}{7.2}\selectfont
    \begin{columns}
        \begin{column}{0.5\textwidth}
            \includegraphics[width=0.8\textwidth]{graphics/rlc_loop.png}
        \end{column}
        \begin{column}{0.5\textwidth}
            \begin{eqnarray*}
            u_{0} & = & \sin\left(t\right)\\
            u_{1} & = & R_{1}i_{1}\\
            u_{2} & = & R_{2}i_{2}\\
            u_{L} & = & L\frac{di_{L}}{dt}\\
            i_{C} & = & C\frac{du_{C}}{dt}\\
            u_{0} & = & u_{1}+u_{2}\\
            u_{L} & = & u_{1}+u_{2}\\
            u_{C} & = & u_{2}\\
            i_{0} & = & i_{1}+i_{L}\\
            i_{1} & = & i_{2}+i_{C}
            \end{eqnarray*}
        \end{column}
    \end{columns}
\end{frame}


\section{Modelica}

\begin{frame}{Modelica}
    \begin{block}<1->{}
        Es un lenguaje de modelado (estándar abierto) que permite la especificación de modelos matemáticos de complejos sistemas naturales o artificiales.
    \end{block}
    Características principales:
    \begin{itemize}
        \item<2-> Modelica es un lenguaje orientado a objetos con un concepto general de clase que unifica las ideas de clases, tipos genéricos y subtipos en un solo lenguaje. Esto facilita la reutilización de componentes y la evolución de los modelos. 
    \end{itemize}
\end{frame}

\begin{frame}{Modelica}
    \begin{itemize}
        \item<1-> Modelica está principalmente basado en ecuaciones en lugar de asignaciones. Esto permite un modelado acausal lo cual facilita la reutilización de las clases ya que las ecuaciones no especifican una dirección particular para el flujo de datos. De esta manera una clase Modelica se puede adaptar a contextos con diferentes flujos de datos.
        \item<2-> Modelica permite describir y conectar componentes de modelos pertenecientes a diferentes dominios como eléctricos, mecánicos, termodinámica, hidráulica, biología, control, etc. 
    \end{itemize}
\end{frame}

\begin{frame}{Conceptos básicos}
    \begin{itemize}
        \item<1-> Un modelo Modelica se construye a partir de clases.
        \item<2-> Los principales elementos que contiene una clase  son declaraciones de variables y ecuaciones.
    \end{itemize}
\end{frame}

\begin{frame}[fragile]{Conceptos básicos}
    \begin{block}{Hello World}
        $\dot{x}(t)=-a x(t)$
    \end{block}
    \pause
    \begin{lstlisting}[language=Modelica]
  class HelloWorld
    Real x(start = 1);
    parameter Real a = 1;
  equation
    der(x) = -a*x;
  end HelloWorld;
    \end{lstlisting}  
\end{frame}
 
\section{Causalización de un modelo DAE}

\begin{frame}{Causalización}
    ¿Como convertimos un modelo DAE acausal en un modelo ODE causal?
    \pause
    \begin{itemize}
        \item<2-> Necesitamos reordenar horizontalmente el modelo de manera que del lado izquierdo de cada ecuación solo aparezca o bien la derivada de una variable de estado o bien una variable algebraica.
        \item<3-> Necesitamos ordenar verticalmente las ecuaciones de manera que puedan ser resueltas secuencialmente.
    \end{itemize}
\end{frame}

\begin{frame}{Causalización}
    \begin{itemize}
        % \item<1-> Definimos que para toda variable $x$ tal que $\dot{x}$ aparece en el modelo, $x$ es conocida y $\dot{x}$ es una incógnita.
        \item<1-> Construimos un grafo bipartito no dirigido a partir del modelo DAE, el cual posee un vértice para cada ecuación del modelo y un vértice por cada incógnita del modelo. Una arista del grafo representa la ocurrencia de una incógnita en una ecuación.
        \item<2-> Calculamos un matching máximo sobre el grafo para definir qué variable vamos a despejar de qué ecuación.
        \item<3-> A partir del grafo bipartito no dirigido y del matching construimos un grafo dirigido.
        \note{
            El grafo dirigido representa las dependencias entre las ecuaciones.
            \begin{itemize}
                \item<1-> Por cada arista del matching creamos un nodo en el grafo dirigido. De esta manera un nodo del grafo dirigido esta asociado a una ecuación y a una incógnita de esa ecuación la cual a sido seleccionada para ser despejada.
                \item<2-> Dados dos nodos, a y b, del grafo dirigido, trazamos una arista de a hacia b si en el grafo original existe una arista, no perteneciente al matching, entre la ecuación asociada al nodo a y la incógnita asociada al nodo b.
            \end{itemize}
        }
        \item<4-> Aplicamos el algoritmo de Tarjan sobre el grafo dirigido.
        \item<5-> Despejamos de cada ecuación la variable seleccionada por el matching.
    \end{itemize}
\end{frame}

\begin{frame}{Algoritmo de Tarjan}
    \begin{itemize}
        \item<1-> Permite encontrar los componentes fuertemente conexos en un grafo dirigido.
        \item<2-> Complejidad $O(\abs{V}+\abs{E})$
        \item<3-> Posee la propiedad de que el orden en el que son identificados los componentes fuertemente conexos constituye un orden topológico reverso del grafo dirigido acíclico formado por los mismos componentes fuertemente conexos.
        \note{
            Ningún componente fuertemente conexo va a ser identificado antes que alguno de sus sucesores.
        }
    \end{itemize}
\end{frame}

\begin{frame}{Algoritmo de Tarjan}
    \framesubtitle{Ejemplo}
    \begin{figure}
        \includegraphics[width=0.6\textwidth]{graphics/tarjan_step_0.png}
    \end{figure}
\end{frame}

\begin{frame}{Algoritmo de Tarjan}
    \framesubtitle{Ejemplo}
    \begin{figure}
        \includegraphics[width=0.7\textwidth]{graphics/tarjan_step_1.png}
    \end{figure}
\end{frame}

\begin{frame}{Algoritmo de Tarjan}
    \framesubtitle{Ejemplo}
    \begin{figure}
        \includegraphics[width=0.7\textwidth]{graphics/tarjan_step_2.png}
    \end{figure}
\end{frame}

\begin{frame}{Algoritmo de Tarjan}
    \framesubtitle{Ejemplo}
    \begin{figure}
        \includegraphics[width=0.7\textwidth]{graphics/tarjan_step_3.png}
    \end{figure}
\end{frame}

\begin{frame}{Algoritmo de Tarjan}
    \framesubtitle{Ejemplo}
    \begin{figure}
        \includegraphics[width=0.7\textwidth]{graphics/tarjan_step_4.png}
    \end{figure}
\end{frame}


\begin{frame}[fragile]{Aplicación del Algoritmo de Tarjan sobre un modelo DAE}
    \begin{columns}
        \begin{column}{0.5\textwidth}
          \centering
          \begin{align*}
            f_1(z_3,z_4) &= 0 \\
            f_2(z_2) &= 0 \\
            f_3(z_2,z_3,z_5) &= 0 \\
            f_4(z_1,z_2) &= 0 \\
            f_5(z_1,z_3,z_5) &= 0 \\
          \end{align*}
        \end{column}
        \begin{column}{0.5\textwidth}
        \begin{figure}
           \centering
           \includegraphics[width=0.8\textwidth]{graphics/bipartito_original.png}
        \end{figure}
        \end{column}
    \end{columns}
\end{frame}

\begin{frame}[fragile]{Aplicación del Algoritmo de Tarjan sobre un modelo DAE}
    \begin{columns}
        \begin{column}{0.5\textwidth}
        \begin{figure}
           \centering
           \includegraphics[width=0.6\textwidth]{graphics/bipartito_matching.png}
        \end{figure}
        \end{column}
        \begin{column}{0.5\textwidth}
        \begin{figure}
           \centering
           \includegraphics[width=0.8\textwidth]{graphics/dirigido.png}
        \end{figure}
        \end{column}
    \end{columns}
\end{frame}

\begin{frame}[fragile]{Aplicación del Algoritmo de Tarjan sobre un modelo DAE}
    \begin{columns}
        \begin{column}{0.5\textwidth}
            \begin{figure}
               \centering
               \includegraphics[width=0.8\textwidth]{graphics/dirigido_componentes.png}
            \end{figure}
        \end{column}  
        \begin{column}{0.5\textwidth}
            \begin{align*}
            z_2 &= f_2'() \\
            z_1 &= f_4'(z_2) \\
            (z_3, z_5) &= f_{35}'(z_2) \\
            z_4 &= f_1'(z_4) \\
            \end{align*}
        \end{column}
    \end{columns}
\end{frame}


\section{Visión General del Compilador ModelicaCC}

\begin{frame}
    La compilación de un modelo Modelica es el proceso mediante el cual se transforma un modelo de alto nivel, el cual puede contener clases y ecuaciones acausales, en un modelo plano, sin clases, y donde las ecuciones se encuentran ordenadas o causalizadas.
    \pause
    El compilador ModelicaCC se planteo como un conjunto de componentes independientes.
    \pause
    \begin{itemize} 
        \item<1-> \textbf{flatter}: Es el componente encargado de el análisis sintáctico y semántico, y del aplanado de la estructura de objetos. 
        \item<2-> \textbf{μModelica}: Es el componente responsable de reescribir el modelo utilizando un subconjunto del lenguaje Modelica denominado μModelica.
        \item<3-> \textbf{antialias}: Es el componente encargado de eliminar ecuaciones triviales de la forma \verb+a=b+.
        \item<4-> \textbf{cauzalize} Es el componente responsable de la causalización del conjunto de ecuaciones del modelo.
\end{frame}

\section{Construcción del grafo bipartito}

\begin{frame}{Representación del grafo bipartito}
    \begin{itemize}
        \item<1-> Utilizamos la biblioteca Boost Graph para la representación del grafo bipartito.
        \item<2-> La BGL ofrece dos clases para representar grafos: adjacency\_list y adjacency\_matrix.
        \item<3-> La adjacency\_list resulta más conveniente que la adjacency\_matrix para nuestro caso. 
        \begin{itemize}
            \item<4-> Mejor complejidad espacial: $O(V+E)$ contra $O(V^{2})$.
            \item<5-> Recorrido de aristas salientes de un vértice más eficiente: $O(V+E)$ contra $O(V^{2})$.
            \item<6-> Más eficiente para el caso de grafos ralos.
        \end{itemize}
    \end{itemize}
\end{frame}

\begin{frame}[fragile]{Definición del grafo}
   \fontsize{9pt}{7.2}\selectfont
    \note{
        La elección de listS como contenedor para las aristas y vértices del grafo se justifica ya que con este tipo de estructura de datos la complejidad temporal de las operaciones que modifican el grafo es mejor. Por ejemplo la operación remove\_vertex(), la cual se utiliza en la optimización del proceso de causalización, es constante para el contenedor listS pero O(V+E) para el contenedor boost::vecS.
    }
    \begin{lstlisting}
struct VertexProperties {
    Modelica::AST::Equation  eq;
    Modelica::AST::Expression unknown;
};

typedef boost::adjacency_list<boost::listS,
        boost::listS,
        boost::undirectedS,
        VertexProperties> CausalizationGraph;

typedef CausalizationGraph::vertex_descriptor Vertex;

typedef CausalizationGraph::edge_descriptor Edge;
    \end{lstlisting}
\end{frame}

\begin{frame}[fragile]{Construcción del grafo a partir de un modelo Modelica}
    A partir del AST y de la tabla de símbolos  del modelo identificamos las ecuaciones, las incógnitas y en qué ecuación aparece cada incógnita.
    \pause
    \begin{itemize}
        \item<2-> Expandimos las ecuaciones de tipo \verb+For-Equation+.
        \begin{lstlisting}[language=Modelica]
            for i in 1:10 loop
                a[i] = i;
            end for;
        \end{lstlisting}
        \item<3-> Idenficamos las incógnitas recorriendo la tabla de símbolos.
        \begin{itemize}
            \item Es una variable continua.
            \item No es un estado.
            \item Es la derivada temporal de una variable.
        \end{itemize}
        \item<4-> "Trazamos" las aristas del grafo determinando en qué ecuación ocurre cada variable.
    \end{itemize}
\end{frame}

\section{Causalización}

\begin{frame}[fragile]{Aplicación del Algoritmo de Tarjan}
    \note{
        El grafo dirigido representa las dependencias entre las ecuaciones.
        \begin{itemize}
            \item<3-> Por cada arista del matching creamos un nodo en el grafo dirigido. De esta manera un nodo del grafo dirigido esta asociado a una ecuación y a una incógnita de esa ecuación la cual a sido seleccionada para ser despejada.
            \item<3-> Dados dos nodos, a y b, del grafo dirigido, trazamos una arista de a hacia b si en el grafo original existe una arista, no perteneciente al matching, entre la ecuación asociada al nodo a y la incógnita asociada al nodo b.
        \end{itemize}
    }
    \begin{itemize}
        \item<1-> Utilizamos la función \verb+checked_edmonds_maximum_cardinality_matching+ provista por la BGL para el cálculo de matching.
        \item<2-> Construimos el grafo dirigido.
        \item<3-> Utilizamos la función \verb+strong_components+ de la BGL.
    \end{itemize}
\end{frame}

\begin{frame}[fragile]{Complejidad}
    \note{
        Matching: $O(V\times E)$ -> $E\geqslant\dfrac{V}{2}$\\
        $O(\dfrac{E}{V})$ expresa la complejidad temporal de las operaciones que recorren la lista de adyacencia de un vértice. Para grafos ralos $\dfrac{E}{V}$ es mucho más chico que $V$ y puede ser considerado constante.\\
        Mapas: vertex2component, matching. Complejidad logaritmica.
    }
    \begin{itemize}
    \item \verb+checked_edmonds_maximum_cardinality_matching+: $O(V^{2})$.
    \item La construcción del grafo dirigido: $O(V\log(V))$.
    \item \verb+strong_components+: $O(V+E)$.
    \end{itemize}
\end{frame}

\section{Optimización del Proceso de Causalización}

\begin{frame}{Un algritmo simple para obtener la causalización de un modelo DAE sin bucles}
    \begin{itemize}
        \item Si una ecuación posee una sola incógnita, entonces debemos utilizar esa ecuación para resolver esa incógnita.
        \item Si una incógnita aparece en una sola ecuación, entonces no tenemos más alternativa que utilizar esa ecuación para resolver esa incógnita.
    \end{itemize}
\end{frame}

\begin{frame}[fragile]{Un algritmo simple para obtener la causalización de un modelo DAE sin bucles}
\begin{columns}
    \begin{column}{0.5\textwidth}
        \begin{align*}
            h_{1}(z_{1},z_{3}) &=0 \\
            h_{2}(z_{2}) &=0 \\
            h_{3}(z_{1},z_{2}) &=0
        \end{align*}
    \end{column}
    \begin{column}{0.5\textwidth}
        \includegraphics[width=0.8\textwidth]{graphics/ej_alg_simple_op_1.png}
    \end{column}
\end{columns}
\end{frame}

\begin{frame}[fragile]{Un algritmo simple para obtener la causalización de un modelo DAE sin bucles}
\begin{columns}
    \begin{column}{0.5\textwidth}
        \includegraphics[width=0.8\textwidth]{graphics/ej_alg_simple_op_2.png}
    \end{column}
    \begin{column}{0.5\textwidth}
        \includegraphics[width=0.8\textwidth]{graphics/ej_alg_simple_op_3.png}
    \end{column}
\end{columns}
\end{frame}

\begin{frame}[fragile]{Un algritmo simple para obtener la causalización de un modelo DAE sin bucles}
\begin{columns}
    \begin{column}{0.5\textwidth}
    \begin{center} 
        \includegraphics[width=0.5\textwidth,height=0.5\textheight]{graphics/ej_alg_simple_op_4.png}
    \end{center}
    \end{column}
    \begin{column}{0.5\textwidth}
        \begin{align*}
            z_{2} &= h_{2}'() \\
            z_{1} &= h_{3}'(z_{2}) \\
            z_{3} &= h_{1}'(z_{1}) 
        \end{align*}
    \end{column}
\end{columns}
\end{frame}

\begin{frame}[fragile]{Complejidad del algoritmo simple}
\begin{center}
    $O(V\log V)$
\end{center}
\end{frame}

\begin{frame}{Algoritmo final de causalización}
\begin{itemize}
    \item<1-> Construimos el grafo bipartito a partir del modelo DAE.
    \item<2-> Aplicamos el algoritmo simple.
    \item<3-> Aplicamos el algoritmo de Tarjan sobre el grafo resultante de aplicar el algoritmo simple.
    \item<4-> Finalmente el despeje de las incógnitas los hacemos utilizando la biblioteca de manipulación simbólica GiNaC.
\end{itemize}
\end{frame}

\section{Ejemplos}

\begin{frame}[fragile]{Circuito RLC}
   \fontsize{8pt}{7.2}\selectfont
    \begin{columns}
        \begin{column}{0.5\textwidth}
        \begin{lstlisting}[language=Modelica]
model RLC
    Real u1; Real u2;
    Real uL; Real iC;
    Real i1; Real i2;
    Real u0; Real i0;
    Real iL; Real uC;
    parameter Real R1;
    parameter Real R2;
    parameter Real L;
    parameter Real C;
    equation
        u0 = sin(time);
        u1 = R1 * i1;
        u2 = R2 * i2;
        uL = L * der(iL);
        iC = C *der(uC);
        u0 = u1 + uC;
        uL = u1 + u2;
        uC = u2;
        i0 = i1 + iL;
        i1 = i2 + iC;
end RLC;
        \end{lstlisting}
        \end{column}
        \begin{column}{0.5\textwidth}
        \begin{lstlisting}[language=Modelica]
model RLC
    Real u1; Real u2;
    Real uL; Real iC;
    Real i1; Real i2;
    Real u0; Real i0;
    Real iL; Real uC;
    parameter Real R1;
    parameter Real R2;
    parameter Real L;
    parameter Real C;
    equation
        u0 = sin(time);
        u2 = uC;
        u1 = (-uC)+u0;
        i2 = u2*(1/(R2));
        i1 = (1/(R1))*u1;
        iC = (-i2)+i1;
        uL = u1+u2;
        der(uC) = iC*(1/(C));
        der(iL) = (1/(L))*uL;
        i0 = i1+iL;
end RLC;
        \end{lstlisting}
        \end{column}
    \end{columns}
\end{frame}

\begin{frame}[fragile]{Circuito RLC con bucle algebraico}
    \begin{columns}
        \begin{column}{0.5\textwidth}
            \includegraphics[width=0.8\textwidth]{graphics/rlc_loop.png}
        \end{column}
        \begin{column}{0.5\textwidth}
            \includegraphics[width=0.8\textwidth]{graphics/rlc_loop_graph.png}
        \end{column}
    \end{columns}
\end{frame}

\begin{frame}[fragile]{Circuito RLC con bucle algebraico}
   \fontsize{7pt}{7.2}\selectfont
   \note{
        En este caso elbucle pudo ser resuelto (todos los lados derechos de las ecuaciones dependen solo de la variable u0) porque las ecuaciones que lo conformaban eran lineales.
   }
    \begin{columns}
        \begin{column}{0.4\textwidth}
        \begin{lstlisting}[language=Modelica]
model RLC
    Real u1; Real u2; Real u3;
    Real uL; Real u0; Real i0;
    Real i1; Real i2; Real i3;
    Real iL;
    parameter Real R1;
    parameter Real R2;
    parameter Real L;
    equation
        u0 = sin(time);
        u1 = R1 * i1;
        u2 = R2 * i2;
        u3 = R3 * i3; 
        uL = L * der(iL); 
        u0 = u1 + u3; 
        uL = u1 + u2;
        u3 = u2;
        i0 = i1 + iL;
        i1 = i2 + i3;
  end RLC;
        \end{lstlisting}
        \end{column}
        \begin{column}{0.6\textwidth}
        \begin{lstlisting}[language=Modelica]
model RLC
    Real u1; Real u2; Real u3;
    Real uL; Real u0; Real i0;
    Real i1; Real i2; Real i3;
    Real iL;
    parameter Real R1;
    parameter Real R2;
    parameter Real L;
    equation
        u0 = sin(time);
        i1 = (1/(R1*R2+R1*R3+R2*R3))*(u0*R3+u0*R2);
        u2 = u0*(1/(R1*R2+R1*R3+R2*R3))*R2*R3;
        i3 = u0*(1/(R1*R2+R1*R3+R2*R3))*R2;
        u1 = (1/(R1*R2+R1*R3+R2*R3))*(u0*R1*R2+u0*R1*R3);
        u3 = u0*(1/(R1*R2+R1*R3+R2*R3))*R2*R3;
        i2 = u0*(1/(R1*R2+R1*R3+R2*R3))*R3;
        uL = u1+u2;
        der(iL) = (1/(L))*uL;
        i0 = i1+iL;
end RLC;
        \end{lstlisting}
        \end{column}
    \end{columns}
\end{frame}

\section{Prueba de Eficiencia}

\begin{frame}{Transmisión de energía en una dimensión}
    \begin{center}
        \includegraphics[width=0.8\textwidth]{graphics/heat_transfer.eps}
    \end{center}
    \pause
    \begin{eqnarray*}
        \dot{u}_{n}(t) & = & 0\\
        \dot{u}_{1}(t) & = & \frac{u_{1}(t)-u_{2}(t)}{\frac{1}{N}}\\
        \dot{u}_{i}(t) & = & \frac{u_{i+1}(t)-u_{i}(t)}{\frac{1}{N}}+\frac{u_{i-1}(t)-u_{i}(t)}{\frac{\gamma}{N}}\;
    \end{eqnarray*}
\end{frame}

\begin{frame}[fragile]{Modelo Modelica}
\fontsize{8pt}{7.2}\selectfont
\begin{lstlisting}[language=Modelica]
  model OneDHeatTransferTI_FD
    parameter Real L = 0.2;
    constant Integer N = 100;
    parameter Real T0 = 273.15;
    parameter Real TN = 330;
    parameter Real cp = 910;
    parameter Real lambda = 237;
    parameter Real rho = 2712;
    final parameter
    Real dx = L / (N - 1);
    Real T[N];
    equation
      der(T[N]) = 0;
      for i in 2:N - 1 loop
        der(T[i]) = lambda * ((T[i + 1] - T[i]) / dx + ((-T[i]) + T[i - 1]) / dx) / cp / rho / dx;
      end for;
      der(T[1]) = lambda * ((T[2] - T[1]) / dx) / cp / rho / dx;
  end OneDHeatTransferTI_FD; 
\end{lstlisting}
\end{frame}

\begin{frame}{Resultados Caso 1}
\begin{center}
\includegraphics[width=0.8\textwidth]{graphics/OneDHeatTransfer.png}
\end{center}
\end{frame}

\begin{frame}[fragile]{Modelo Modelica con Bucle Algebraico}
\fontsize{8pt}{7.2}\selectfont
\begin{lstlisting}[language=Modelica]
  model OneDHeatTransferTI_FD
    ...
    Real T[N];
    Real U[N]; 
    equation
      der(T[N]) = 0;
      U[1]*U[2]=0;
      for i in 2:N - 1 loop
        der(T[i]) = der(T[i-1]) + lambda * ((T[i + 1] - T[i]) / dx + ((-T[i]) + T[i - 1]) / dx) / cp / rho / dx;
        U[i] = U[i+1]+U[i-1]; 
      end for;
      der(T[1]) = lambda * ((T[2] - T[1]) / dx) / cp / rho / dx;
      U[N] = U[N-1]*2+sin(U[N]); 
  end OneDHeatTransferTI_FD; 
\end{lstlisting}
\end{frame}

\begin{frame}{Resultados Caso 2}
\begin{center}
\includegraphics[width=0.8\textwidth]{graphics/OneDHeatTransfer_loop.png}
\end{center}
\end{frame}

\section{Conclusiones}

\begin{frame}{Conclusiones}
    \begin{itemize}
        \item<1-> Presentamos la implementación en C++ de la etapa de causalización de eucaciones de un compilador modelica.
        \item<2-> Estudiamos e implementamos tres estrategias distintas para resolver el problema de causalización analizando la complejidad computacional de cada una de ellas.
        \item<3-> Utilizamos bibliotecas abiertas como la BGL o GiNaC.
        \item<4-> Utilizando ejemplos del dominio de la ingenería electrica probamos que algoritmo de causalización produce resultados correctos para casos con y sin bucles algebraicos.
        \item<5-> Analizamos el desempeño de las diferentes estrategías utilizando un modelo de grande.
    \end{itemize}
\end{frame}

\begin{frame}[fragile]{Trabajo a futuro}
    \note{
        Singularidad estructural: Un modelo presenta singularidades estructurales si el orden del sistema
    }
    \begin{itemize}
        \item Implementación del algoritmo de Pantelides para eliminar las singularidades estructurales.
        \item Tratamiento de los lazos algebraicos mediante la técnica de Tearing.
        \item Causalización de ecuaciones de tipo \verb+For-Equation+ sin realizar la expansión del For.
    \end{itemize}
\end{frame}

\begin{frame}{}
\begin{center}
    ¿Preguntas?
\end{center}
\end{frame}


\end{document}
